\documentclass[11pt]{article}
\usepackage[margin=1in]{geometry}
\usepackage{fontspec}
\usepackage{xeCJK}
\IfFontExistsTF{Times New Roman}{\setmainfont{Times New Roman}}{\setmainfont{TeX Gyre Termes}}
\IfFontExistsTF{Songti TC}{\setCJKmainfont{Songti TC}}{\setCJKmainfont{Noto Serif CJK TC}}
\usepackage{amsmath}
\usepackage{graphicx}
\usepackage{booktabs}
\usepackage{tabularx}
\usepackage{array}
\usepackage{float}
\usepackage{placeins}
\usepackage[hidelinks]{hyperref}
\usepackage{microtype}
\setlength{\parindent}{1.5em}
\setlength{\parskip}{0.25em}
\title{Refusing Self-Declared Relationships Before Prompting Mitigates Unauthenticated External Traffic in LLM-Agent Collectives}
\author{Ben Tseng\\\small Draft for feedback, 11 November 2026}
\date{}
\begin{document}
\maketitle


\begin{abstract}

Multi-agent LLM systems let agents pool reasoning, but they also let outsiders manufacture apparent social support. We show that how a router handles one piece of metadata, the relationship a sender claims to the receiver, decides whether that support can be bought. In an eight-agent collective built on the \textit{Physics of Agents} substrate, each added external sender delivers the same false answer while claiming a positive relationship. Across two official deployments (Qwen3.5-27B and Mistral Small 4; 3,600 episodes in 20 complete model-seed blocks), trusting the claim made each additional sender multiply the conditional odds of an incorrect majority by 1.90 (95\% Wald CI 1.59 to 2.27). Resolving the claim before prompt construction changed this. Downgrading unverified traffic to a labeled unverified-source channel, while keeping it readable, roughly halved the odds (OR 0.48, 0.32 to 0.71), and dropping it returned the collective to control-level behavior. Both effects replicated in each deployment, held under within-question randomization tests, and kept their direction under the preregistered pooled model (OR 1.23 and 0.79). Relationship claims should therefore be treated as authorization data resolved outside the model: a deterministic, auditable routing rule protects collective decisions without judging message truth. These results cover external attackers without credentials; testing compromised members and token forgery or replay is left to future work.

\textbf{Keywords:} multi-agent security, unauthenticated traffic, message authentication, receiver-side routing, collective decision-making, trustworthy AI

\end{abstract}
\section{Introduction}

Large language model (LLM) agents increasingly solve tasks through structured conversation rather than isolated inference. Role-playing systems, multi-agent debate, and programmable conversation frameworks let agents divide work, challenge intermediate answers, and combine distinct traces \cite{guo2024survey,li2023camel,wu2024autogen,du2024debate,liang2024divergent}. Their benefits depend on the communication graph and the semantics attached to its edges: sparse interaction can preserve performance while reducing cost, so routing is part of the computation rather than neutral transport \cite{li2024sparse}. This paper examines one narrow security question at that boundary: what happens when external peer traffic is allowed to choose its own apparent relationship to the receiver?

The challenge is not only malicious language; it is preserving provenance and authorization while messages move between agents. Routers see message content and sometimes sender, role, or relationship metadata, but not the process that produced them; whatever metadata the router accepts then shapes how the receiving model is prompted. LLM collectives respond to numerical majorities, apparent expertise, speaking dominance, and persuasive rhetoric \cite{chuang2024opinion,choi2025conformity,ashery2025conventions,zomer2026collective,ko2026social,oh2026belief,amayuelas2024attack,kraidia2026persuasion}, and repeated content alone can drive much of this conformity even when no speaker is named \cite{hu2026speaker}. Communication attacks can alter inter-agent messages \cite{he2025redteaming}, prompt infections can propagate between agents \cite{lee2024promptinfection}, conformity-driven injections can push a debate toward a wrong consensus \cite{cui2025madspear}, and delegation can erase source and authority metadata while retaining actionable content \cite{zhou2026trustcascade}. These findings make prompt construction a trust boundary: a router can amplify unverified traffic before the model evaluates its substance.

Existing defenses already treat provenance as a first-class signal. Spotlighting marks untrusted input so the model can tell sources apart \cite{hines2024spotlighting}; instruction hierarchies train models to privilege trusted instructions \cite{wallace2024hierarchy}; capability- and information-flow-based designs keep untrusted data from steering control flow \cite{debenedetti2025camel}, \cite{costa2025fides}; and recent multi-agent work signs inter-agent messages and quarantines unsigned ones \cite{paul2026beyond}. Topology-guided systems detect anomalous utterance patterns and intervene in a communication graph \cite{wang2025gsafeguard}, and Byzantine- or credibility-weighted aggregation limits the influence of faulty agents \cite{zheng2025byzantine}, \cite{ebrahimi2025credibility}. These approaches establish that isolating unverified input is useful. They do not ask what happens when the unverified element is a \textit{relationship claim} that decides which social channel a message enters, nor how such a claim changes a collective's terminal decision. Authentication also cannot determine whether a valid sender is truthful. The question addressed here is therefore narrower: if a router keeps relationship authorization separate from semantic judgment, does that separation change collective outcomes?

The experiment operationalizes this problem as a bundled external-traffic treatment. Each additional nominal external sender contributes the same fixed false-option message and targets exactly four honest agents. The receiving prompt renders each message as a bare bullet of text and does not display sender identifiers, so additional senders reach the model only as additional identical copies of the false message and as wider reach across the collective. The budget therefore manipulates repetition and reach, not the number of apparent speakers. This matches a practical abuse path—one principal creates more identifiers to deliver more apparent corroboration—but the model never observes the identifiers themselves. Throughout the paper, causal claims refer to this repetition-and-reach bundle, not to an identity-count mechanism.

We ask: \textbf{How does the bundled budget of unauthenticated external senders and standardized false messages affect incorrect-majority outcomes, and to what extent do credential-aware downgrade and drop policies mitigate the effect across two official LLM deployments?} The router verifies an HMAC token issued over the sender, receiver, signed relationship, and an episode-step label. Enrollment and key distribution are assumed, while truthfulness, key compromise, and general identity proofing remain outside scope. Invalid traffic is either trusted as claimed, retained in an unverified-source channel, or dropped before prompt construction. In this experiment every external envelope lacks a token and comes from an identifier outside the enrolled graph, so the credential check and a plain membership check ("is this sender an enrolled node?") route traffic identically. The data therefore test whether a router should honor self-declared relationships from non-members; they do not test the additional receiver, relation, or freshness bindings of the token, which are design rationale rather than evaluated properties.

We test two confirmatory hypotheses. H1 predicts that, without enforcement, adding external senders increases the odds of an incorrect terminal majority. H2, the preregistered primary contrast, predicts that routing unverified traffic to a labeled unverified-source channel reduces those odds relative to trusting the claim. A third contrast, \texttt{drop} versus trust, serves as a manipulation check: because \texttt{drop} removes external traffic before prompt construction, it should return the collective to control behavior if the attack works only through the routed channel. An exploratory analysis (H4) connects the policies to the critical-budget language of collective dynamics.

This paper makes three contributions. First, it identifies relationship claims as an attack surface distinct from message content. In signed-graph collectives, the claimed relationship dictates the social channel a message enters. Therefore, an outsider who controls this claim controls the standing their message receives from the collective. Second, it shows that a deterministic receiver-side rule that resolves the claim before prompting protects collective decisions in two official deployments, and that the two policies differ in a useful way: \texttt{downgrade} keeps unverified information available yet still halves the attack's effect, while \texttt{drop} removes it entirely. Third, it provides a controlled, reproducible protocol on the \textit{Physics of Agents} substrate \cite{el2026physics}: two deployments, 3,600 episodes in complete model-seed blocks, technical admission gates, and a full record of deviations from the registration (Section 3.7), which later work can extend to stronger attackers.

The main findings are clear. Trusting self-declared relationships made the collective steadily more vulnerable: each added sender multiplied the conditional odds of an incorrect majority by 1.90 in the pooled model. \texttt{Downgrade} cut those odds by about half (OR 0.48), and \texttt{drop} returned behavior to control levels. Both effects held in each deployment and under within-question randomization tests, and the preregistered pooled specification gave the same directions. The evidence concerns the admission of self-declared relationships from non-members; Section 6 discusses what stronger attackers, broader tasks, and larger question sets would add.

\section{Related Work}

\subsection{Collective reasoning and communication topology}

LLM multi-agent research treats interaction as a means of expanding the effective computation available to a task. CAMEL showed that role-conditioned agents can sustain autonomous cooperation through conversation \cite{li2023camel}. AutoGen generalized this pattern into programmable multi-agent workflows with configurable participants and termination rules \cite{wu2024autogen}. Multi-agent debate has improved factuality and reasoning on several benchmarks by allowing agents to propose, criticize, and revise answers \cite{du2024debate}. Debate can also encourage divergent reasoning when a controlled protocol prevents early convergence \cite{liang2024divergent}.

These gains are conditional. Social-psychological traits and interaction patterns affect conformity, consensus, and efficiency \cite{zhang2024social}. Sparse communication topologies can match or outperform all-to-all debate \cite{li2024sparse}, showing that more messages are not automatically better. Strong single-agent prompting can approach the best discussion methods in some settings, so collaboration should not be treated as a universal performance multiplier \cite{wang2024bounds}. The relevant unit of analysis is the joint system: model, prompt, topology, update rule, and aggregation procedure.

The focal \textit{Physics of Agents} work models interacting LLM agents on signed graphs as a compact statistical-mechanics system with indifference, polarization, and consensus regimes, and fits coupling parameters—separate weights for agreeing, disagreeing, and connection-level ties—that predict individual opinion updates \cite{el2026physics}. The present experiment reuses its substrate (Section 3.2). Related network studies likewise find that information flow and cognitive update rules shape collective behavior \cite{chuang2024opinion}, \cite{zomer2026collective}. That work contains no external adversary or security intervention. A credential decision changes who may exert which kind of influence; in the coupling language of \cite{el2026physics}, \texttt{downgrade} plausibly moves an external edge from the agreeing coupling toward the connection-level coupling. Testing that mapping quantitatively requires deployment-matched coupling estimates and observations near the relevant threshold.

\subsection{Conformity, consensus, and collective failure}

Collective reasoning fails when interaction amplifies shared mistakes. Networked LLM agents often move toward truth-aligned consensus under neutral conditions, but prompted bias can alter their trajectories \cite{chuang2024opinion}. Group-conformity experiments show that neutral agents respond to numerical majorities and that apparently capable agents may exert disproportionate influence \cite{choi2025conformity}. These findings make sender count a potential control variable even when the message content is simple.

The present study differs from a generic conformity experiment in one respect: the manipulated senders lack valid membership credentials. Their messages are therefore not simply another social condition. Placing unverified copies of a claim in the trusted-relationship section of a prompt is a failure in the representation of provenance: the router, not the model, has already granted the content the standing of a trusted neighbor. Because repetition without any named speaker already drives much LLM conformity \cite{hu2026speaker}, the relevant security variable is how many copies of unverified content reach the trusted channel, not how many identities the model perceives.

\subsection{Adversarial collaboration and external-input attacks}

Research on collaboration attacks shows that one malicious debate participant can reduce group accuracy, especially when it generates persuasive arguments \cite{amayuelas2024attack}. A single persuasive adversary can also raise consensus on incorrect answers, and adding agents or rounds does not reliably mitigate the effect \cite{kraidia2026persuasion}. MAD-Spear compromises a few debaters so that they emit several plausible but incorrect responses, exploiting conformity to make a wrong answer look independently corroborated \cite{cui2025madspear}. These studies establish that collaborative protocols must tolerate adversarial speech, but they usually begin with an adversary that already occupies a recognized participant slot. An impersonation embedded inside a validly credentialed message, as in MAD-Spear, would pass the check studied here.

More direct communication attacks manipulate messages in transit rather than compromise an enrolled agent. Agent-in-the-Middle intercepts inter-agent messages and generates context-aware malicious instructions within a role-restricted communication setting \cite{he2025redteaming}. Prompt Infection shows that injected prompts can propagate from one agent to another and evaluates "LLM Tagging," which marks agent-origin text; tagging alone gave little protection but helped in combination with other defenses \cite{lee2024promptinfection}. Paul and Nandy sign every inter-agent message with a provenance record and quarantine unsigned messages, reducing inter-agent injection success in a six-agent pipeline \cite{paul2026beyond}. G-Safeguard detects anomalies in an utterance graph and changes topology to limit spread \cite{wang2025gsafeguard}. The present intervention shares the verify-then-isolate principle of \cite{paul2026beyond} but differs in what is verified and what is measured: the credential covers a relationship-conditioned route rather than only sender identity, the \texttt{downgrade} policy keeps content while removing relationship privilege, and the outcome is a collective decision on a signed graph rather than single-agent injection success.

Tool-integrated agents face a related external-input problem. InjecAgent demonstrates that indirect prompt injection embedded in retrieved content can cause harmful actions \cite{zhan2024injecagent}. AgentDojo evaluates attacks in a dynamic task environment and highlights the security-utility tradeoff when defenses reject untrusted content \cite{debenedetti2024agentdojo}. Agent Security Bench broadens the evaluation across attack classes, tools, defenses, and model backbones \cite{zhang2025asb}. Structural defenses mark or isolate untrusted input rather than classify it: spotlighting transforms untrusted text to signal its provenance \cite{hines2024spotlighting}, instruction hierarchies train models to subordinate lower-privilege inputs \cite{wallace2024hierarchy}, and CaMeL and FIDES separate control flow from untrusted data with capabilities or information-flow labels \cite{debenedetti2025camel}, \cite{costa2025fides}. Here the untrusted object is a peer message plus relationship metadata, not a retrieved tool result. The router does not classify the prose as malicious; it applies a fixed verification predicate. The \texttt{downgrade} policy is close in spirit to spotlighting, because it labels the provenance of retained content in the prompt; Section 5.4 discusses what that implies for interpreting H2.

Provenance can also be lost during legitimate delegation. Trust Cascade reports that inter-agent delegation strips source and authority markers while often preserving actionable payloads \cite{zhou2026trustcascade}. That failure motivates preserving trust metadata across hops. Our experiment addresses a complementary case at ingress: an external sender presents relationship metadata that was never valid. Both cases imply that the receiving prompt should not be the first component asked to recover provenance from prose.

\subsection{Identity, credentials, and deterministic enforcement}

The Sybil attack formalizes the danger of assuming that distinct identifiers represent distinct entities \cite{douceur2002sybil}. Without a trusted certification mechanism or a scarce resource that binds identity to an entity, redundancy can be illusory. Practical Byzantine fault tolerance reaches a complementary conclusion from the protocol side: resilience guarantees depend on a defined replica set and bounded faulty membership \cite{castro1999pbft}. LLM-specific work applies these ideas through confidence-weighted Byzantine-robust consensus \cite{zheng2025byzantine} and learned credibility scores that discount adversarial agents \cite{ebrahimi2025credibility}. These aggregation defenses operate after messages are admitted; the present study concerns the earlier admission decision. Before counting votes or messages, the system must know what constitutes a participant.

Multi-agent security surveys have long separated communication protection, access control, trust, and reputation from application-level reasoning \cite{jung2012security}. Recent work on authenticated delegation makes this separation explicit for AI agents by calling for auditable authority chains and scoped permissions \cite{south2025delegation}. Implementations based on decentralized identifiers and verifiable credentials likewise place identity proof outside the LLM and question how far security procedures can be left to probabilistic text generation \cite{garzon2026credentials}. HMAC supplies a standardized mechanism for source authentication and message integrity when communicating parties share keying material \cite{nist2008hmac}.

Authentication, enrollment, relationship authorization, and routing are distinct layers. The experimental HMAC token is issued over a sender, receiver, relation, and episode-step label. It presupposes that the eight honest agents were enrolled through the experimental issuer. The verified relation then authorizes a route. None of these steps proves that the content is true or that a key is uncompromised. This decomposition prevents the paper from treating cryptographic validity as a complete trust judgment.

\section{Methods}

\subsection{Confirmatory design and deployment boundaries}

We conducted a two-deployment confirmatory experiment under a frozen design. The original preregistration named different deployments; Section 3.7 documents how the final pair was reached. The accepted deployments were Qwen3.5-27B through Alibaba Cloud Model Studio's official Singapore endpoint and the version-pinned \texttt{mistral-small-2603} deployment through Mistral AI's official API. Provider infrastructure, serving configuration, and model family were not controlled as separate factors. We therefore treat the deployments as independent replications and report their estimates separately before the pooled model.

The confirmatory corpus contained 18 held-out, text-only mathematics questions. Two candidate questions that required Asymptote diagrams were excluded before outcome collection. Ground-truth labels were balanced across the two answer options, with nine A and nine B answers. Each deployment was evaluated on ten graph and attack seeds. For each seed, the design crossed attack budgets of one, two, and four external identities with three routing policies: \texttt{none}, \texttt{downgrade}, and \texttt{drop}. A budget-zero control completed the ten-cell block. Each model-seed block therefore contained 180 episodes, and the 20 accepted blocks contained 3,600 episodes.


\begin{table}[htbp]
\centering
\caption{Frozen experimental design and accepted sample}
\scriptsize
\resizebox{\textwidth}{!}{%
\begin{tabular}{lr}
\toprule
Component & Frozen value \\
\midrule
Deployments & Qwen3.5-27B; Mistral Small 4 (\texttt{mistral-small-2603}) \\
Held-out questions & 18 text-only mathematics questions; labels balanced 9/9 \\
Honest agents & 8 \\
Nonzero graph entries & 28, corresponding to 14 undirected edges \\
Interaction & 2 synchronous updates; 3 opinion samples per agent and state \\
External-identity budgets & 0, 1, 2, 4 \\
Recipients per external sender & 4 honest agents \\
Routing policies & none, downgrade, drop \\
Seeds & 10 per deployment \\
Accepted blocks and episodes & 20 blocks; 3,600 episodes \\
\bottomrule
\end{tabular}%
}
\end{table}
\FloatBarrier

\subsection{Honest-agent network and episode procedure}

The honest-agent substrate is inherited from the \textit{Physics of Agents} codebase \cite{el2026physics}, of which this study's repository is a fork. Inherited components are the binary-choice recasting of MATH problems, the personas (each embeds one representative problem the agent has already solved together with its ground-truth solution), the signed interaction graphs with entries in \{−1, 0, +1\}, repeated opinion sampling, and the relationship-conditioned inboxes ("colleagues you tend to agree with" and "colleagues you tend to disagree with"). This study adds the external senders, the credential authority, the three routing policies, the unverified-source inbox, the technical gates, and the analysis.

Each episode instantiated eight honest agents on a signed, symmetric interaction graph. The adjacency matrix had 28 nonzero entries, which corresponded to 14 undirected edges. Edge signs represented positive or negative relationships and determined the relationship-conditioned inbox used by the opinion-update prompt. The graph was fixed within a seed and varied across seeds according to the frozen generator; seed s used the same graph, personas, and attack targets in both deployments and in every cell.

Agents answered one of the held-out binary-choice mathematics questions. The opinion prompt, inherited from \cite{el2026physics}, asks for a quick judgment rather than a worked solution ("Do NOT work through the problem … snap judgement"), which likely increases sensitivity to social cues relative to agents that reason at length. An episode began with independent opinion sampling, followed by two synchronous communication and update steps. At each state, three opinion samples were drawn per honest agent, and the agent's state was the majority sign of the three samples, with ties broken by the first sample. The runner converted parsable answers into signed states and retained a raw zero when a spin call could not yield a valid answer. Messages for a step were collected before any recipient updated, preventing within-step ordering from granting an early sender additional influence.

The terminal outcome was computed only from the eight honest agents. External identities supplied messages but did not vote in the final population. This choice prevented a trivial attack in which the attacker wins by adding its own votes directly to the count. A successful episode instead required external traffic to change enough honest states that the honest-agent majority became incorrect.

\subsection{Bundled external-traffic threat model}

Every external sender transmitted the same deterministic two-sentence template: \textit{I favor option [false letter]: [false choice text]. The available evidence strongly supports option [false letter].} It claimed a positive relationship and carried no credential. Each nominal sender selected four distinct honest recipients through a deterministic episode-seeded procedure; the same target set was reused in both interaction steps. At budget b, the treatment therefore added exactly 4b external envelope deliveries per step and 8b across the episode.

This exposure rule is the central identification boundary. The receiving prompt renders every message, honest or external, as a bare bullet of text without a sender identifier, so the model never observes how many distinct senders produced the copies it reads. From the model's side, raising the budget increases (a) the number of identical false-message copies in a targeted receiver's prompt—at most 1, about 2, and about 3.5 copies per receiver per step at b = 1, 2, and 4—and (b) the number of honest agents directly reached—about 4, 6, and 7.5 of 8 on average. The budget coefficient therefore estimates the effect of repetition plus reach. An identity-count mechanism is not merely unidentified; it is excluded by the prompt format. Testing one would require rendering sender identifiers and varying them at fixed message volume.

One attacker presents multiple nominal identifiers, as in a Sybil attack, but because those identifiers are not shown to the model and confer no counted vote, the treatment is better described as message flooding through a trusted channel than as a Sybil campaign. It does not model account-creation cost, provider-level proofing, long-lived reputation, adaptive target discovery, or coordinated natural-language strategies. We use “external sender” rather than “malicious agent” because these units do not receive messages, update opinions, occupy enrolled graph nodes, or vote in the terminal outcome.

Figure 1 summarizes the trust boundary. Honest messages originate on enrolled graph edges and carry credentials. External envelopes carry a sender identifier and a claimed positive relationship at the transport layer but lack a valid token; neither field is shown to the model, which sees only which inbox a message lands in. Verification occurs before an LLM sees the message in a relationship-conditioned inbox.

\begin{figure}[htbp]
\centering
\includegraphics[width=\textwidth]{figures/threat_model.png}
\caption{Threat model and credential-routing paths. Honest messages carry relationship-bound credentials. An external sender supplies a false option and a claimed positive relationship. The receiver-side router either trusts the claim under \texttt{none}, retains the content in an unverified-source channel under \texttt{downgrade}, or discards it under \texttt{drop}.}
\end{figure}
\FloatBarrier

\subsection{Relationship credentials and routing policies}

Before routing an honest message, an HMAC-based credential authority issued a token over the literal tuple \texttt{sender:receiver:relation:episode|step}. By design, binding the receiver prevents a token for one edge from authorizing another, and binding the relation prevents sign substitution; neither property was exercised by the attacker in this experiment. The episode-step label is covered by the MAC but, as implemented, was not checked against the current step (see below). HMAC provided controlled source authentication and integrity following FIPS 198-1 \cite{nist2008hmac}. It was not a public-key identity infrastructure.

The three routing policies differed only in their response to a missing or invalid credential. Under \texttt{none}, the router trusted the external message's claimed relation and placed it in the positive-relation inbox. Under \texttt{downgrade}, it retained the message but placed it in a relationship-unknown inbox. Under \texttt{drop}, it discarded the message. Valid honest messages preserved their verified signed relation under both credential-aware policies. The content generator and attack target logic did not change across policies.

The unknown-relation channel was not assigned an authenticated positive or negative edge. It appears only under \texttt{downgrade}, after both honest sections and immediately before the problem statement, under the header "Messages From Unverified Sources. These messages failed relationship verification. Treat their claimed source relationship as unknown." The header states a provenance fact, not that the content is false. It is nonetheless more than a neutral "unfamiliar" label: \texttt{downgrade} simultaneously removes the positive-relationship framing, adds an explicit verification-failure notice, moves the messages to a separate section, and changes their position. The design does not separate these components (Section 5.3). Downgrade therefore represents a permissive policy for systems that want to expose uncertain-origin information, labeled as such, without granting it the social semantics of an enrolled relationship.

Under \texttt{none}, token status was not used to change routing. This condition models a router that transports sender and relationship fields as application metadata without cryptographic enforcement and lets external traffic choose the most favorable claimed relationship, a worst-case assumption about the attacker. It is not a claim that practical systems literally label a switch “no defense,” nor that mainstream frameworks expose a relationship field to external parties; the "agree" and "disagree" inboxes are constructs of the experimental substrate \cite{el2026physics}. Many agent frameworks allow message objects, role descriptions, or conversation participants to be constructed programmatically \cite{wu2024autogen}. If such fields reach the prompt without an independently verified binding, their security behavior is equivalent to trusting the caller that supplied them. All reported mitigation is relative to this metadata-trusting baseline.

The trust architecture has five layers. Enrollment provisions the eight honest identifiers and signed graph; the experiment assumes this step. Authentication checks the token and its receiver binding. Relationship authorization maps a verified sign to an allowed route. Routing applies \texttt{none}, \texttt{downgrade}, or \texttt{drop} before prompt construction. Deliberation then asks the LLM to solve the task using the selected channels. A valid token does not establish truthfulness, and the design does not test stolen keys, issuer compromise, revocation, replay, or authenticated malicious agents.

Two implementation limits bound these layers. First, the verifier recomputes the MAC over the episode-step label carried in the token but does not compare it with the current episode and step, so the label is integrity-protected but does not enforce freshness; a captured honest token could be replayed within a seed. Second, the HMAC key for each seed is derived deterministically from a public string so that runs are reproducible. Neither limit affects the reported results, because no external sender ever presented a token, but both mean that the implementation demonstrates the routing decision rather than a hardened credential system. Because every external envelope was untokened and came from a non-member identifier, a membership check on the sender identifier would have produced identical routing.

\subsection{Outcomes, admission gates, and block replacement}

The primary outcome was an indicator that the terminal majority of honest agents was incorrect. A tied terminal population was not counted as attack success. Secondary outcomes included the truth-aligned terminal margin, fraction of correct agents, incorrect consensus, raw zero spins, request errors, exhausted parse failures, recovered parse retries, unique-sender message-call failures, invalid honest credentials, and the fitted 50\% critical budget.

Collection proceeded in complete model-seed blocks. Eligibility required all 180 episode files and no invalid honest credentials. Raw zero spins could not exceed 1\% of spin calls, and unique-sender message-call failures could not exceed 1\% of message calls. These criteria were technical gates, not outcome-dependent filters. A failing block was excluded in full rather than repaired by retaining successful cells.

Five blocks were replaced under dated recollection amendments, each written after the block's gate failure and before its replacement was requested: Qwen seeds 0 and 4, and Mistral seeds 1, 7, and 9. The original preregistration specified only that collection pause for diagnosis after a gate failure; the full-block recollection rule was adopted in the 2026-10-29 amendment and applied to every subsequent failure. Each replacement used a new output directory and reran all ten cells. No predecessor episode entered the accepted manifest. The analysis script enumerated the 20 accepted source files explicitly and locked their SHA-256 hashes before merging.

The replacement rule was applied mechanically: all five failed blocks were recollected, each replacement passed on its first attempt, and no failed block was retained. Per-episode outcome tables for failed blocks were generated automatically before each amendment, so the rule cannot be shown to be blind to outcomes; Section 4.5 reports how H2 changes if the failed attempts are used instead. The rule prevented an analyst from retaining successful cells and rerunning only technically poor cells, a practice that could couple admission to unobserved provider conditions or to outcome inspection. The accepted-block manifest records the source path for every model-seed pair and identifies the replacement directory where applicable. Counts reported in the paper were recomputed from those explicit sources rather than from a recursive search of the data tree.

\subsection{Statistical analysis}

The preregistration specified a mixed-effects logistic model for the incorrect-terminal-majority indicator with fixed effects for bundled attack budget, defense, deployment, and the budget-by-defense interaction, and random intercepts for question and seed/graph. The estimator details below (linear raw budget coding, control coded as \texttt{defense=none}, the 7/3 contrast point, MAP estimation, and priors) were fixed in the analysis code before the confirmatory data were merged but were not written into the preregistration (Section 3.7). In compact form,

\[
\operatorname{logit}\Pr(Y_{iqms}=1)=\beta_0+\beta_1 B_i+\beta_2 D_i+\beta_3 M_i+\beta_4(B_i\times D_i)+u_{q(m)}+v_s,
\]

where \(B_i\) is the raw, unstandardized bundled budget in \(\{0,1,2,4\}\), \(D_i\) is routing policy, \(M_i\) is deployment, \(u_{q(m)}\) is a random intercept for question \(q\) within deployment \(m\), and \(v_s\) is a seed random intercept shared across deployments, which used identical graphs and targets for a given seed. The reference policy was \texttt{none}; the single budget-zero control in each seed was coded as \texttt{defense=none}. Maximum a posteriori fits used \texttt{statsmodels} 0.14.4 \texttt{BinomialBayesMixedGLM.fit\_map}, BFGS with at most 2,000 iterations, a zero-mean Gaussian fixed-effect prior with SD 5.0, and a zero-mean Gaussian prior with SD 0.5 on each log random-effect SD. Fixed-effect covariance was the fixed-effect block of the Laplace/MAP covariance matrix returned by \texttt{cov\_params()}. CIs and p values for coefficients and linear contrasts are two-sided normal-Wald quantities, \(\hat\theta\pm1.96\operatorname{SE}(\hat\theta)\) and \(2\Phi(-|\hat\theta/\operatorname{SE}(\hat\theta)|)\), exponentiated for OR reporting. They are approximate MAP/Laplace intervals rather than likelihood-profile or fully Bayesian credible intervals. No optimizer warning was recorded for any reported fit.

\textbf{Pooled specification.} The preregistered specification used one question intercept \(u_q\) shared by both deployments. Question difficulty, however, is negatively correlated across the two deployments (correlation of clipped logit question rates −0.49): several questions that Qwen always failed were always solved by Mistral, and vice versa. A shared intercept cannot represent this, and it pulls the pooled odds ratios outside the range of both deployment-specific estimates. After review we therefore report the question-within-deployment model above as the primary pooled specification and report the preregistered shared-intercept specification alongside it. Both give the same directions; the shared-intercept model gives smaller magnitudes (Table 4). This change is listed as a deviation in Section 3.7.

H1 estimated \(\exp(\beta_1)\), the subject-specific conditional OR per one-unit increase in the bundled budget under \texttt{none}, holding question and seed random intercepts fixed. H2 compared \texttt{downgrade} with \texttt{none} at the mean of the positive tested budgets, \((1+2+4)/3=7/3\). H3 used the same contrast definition for \texttt{drop}. These are also conditional contrasts. Deployment-specific models used the same defense and raw-budget structure and were interpreted as within-deployment estimates; they were not separately preregistered. Because logistic ORs are non-collapsible, conditional OR magnitudes need not match ORs computed from marginal cell proportions.

\textbf{Near-separation and prior dependence.} Most questions produced the same outcome in every episode of a deployment. Only 4 of 18 Qwen questions and 10 of 18 Mistral questions had any variation in the outcome (Section 4.3). For the remaining questions the random intercept is held back only by its prior, so the estimated question-level random-effect SDs, and therefore the magnitude of the conditional ORs, depend on the prior on the log random-effect SD: in the Qwen-only model the question SD ranges from 1.7 to 4.2 and the H3 OR from 0.13 to 0.06 as that prior SD moves from 0.1 to 0.5, while signs and significance are unchanged. Looser priors do not yield a stable joint MAP maximum. We therefore interpret directions and significance, not conditional OR magnitudes, and treat the random-effect SDs as prior-dependent rather than as estimates. As a consistency check, conditional fitted cell probabilities averaged over posterior random effects reproduced the 20 observed cell rates within 0.0055 (Qwen) and 0.026 (Mistral); this check can detect omitted fixed-effect terms but not random-effect misspecification.

\textbf{Randomization inference.} Because policies were randomized within each question, seed, and budget, with graphs and targets held fixed, we also report exact-style within-question randomization tests of \texttt{none} versus \texttt{downgrade} and \texttt{none} versus \texttt{drop} for each deployment. Policy labels were swapped within each question × seed × budget pair (sign-flip of paired differences, 20,000 draws, seed 20261111). These tests condition on the 18 questions and do not support inference about new questions.

Nine secondary pairwise policy contrasts at budgets one, two, and four were adjusted by the Holm procedure. Each is a point on the fitted linear budget-by-defense line rather than a test of an individual cell. Critical budgets were defined as the budget at which a nondecreasing logistic attack-success curve crossed 50\%. Curves were estimated separately by deployment and policy from observed budgets 0, 1, 2, and 4; when the unconstrained slope was negative, it was set to zero, which makes the crossing infinite by construction. Confidence intervals used 2,000 question-cluster bootstrap resamples with random seed 20261106. Bootstrap samples whose fitted slope was zero or negative were retained as infinite values, and ratio samples that were undefined (both crossings infinite, or a \texttt{none} crossing at or below zero) were counted as infinite. The downgrade-to-none critical-budget multiplier supplied the directional H4 quantity. Because no observed condition approached a 50\% crossing within budget four, all critical budgets were treated as extrapolations.

Post-review sensitivity analyses tested assumptions without changing the hypotheses. First, the pooled MAP model was refit over a 3 x 3 prior grid: fixed-effect prior SDs 2.5, 5, and 10 crossed with log-random-effect-SD prior SDs 0.1, 0.25, and 0.5, under both pooled specifications; the deployment-specific models were refit over the same three random-effect priors. Second, a categorical ten-cell binomial generalized estimating equation included deployment and seed fixed effects, an exchangeable working correlation, and sandwich covariance clustered over the 18 questions. Its H1 checks compared each \texttt{none} budget with control; H2 and H3 averaged the three positive-budget log-odds differences. Third, the pooled MAP model was refit 18 times, omitting one question at a time, and the Qwen model was refit omitting each pair of its four informative questions. Fourth, the pooled model was refit with the five failed first attempts substituted for their accepted replacements. These analyses are robustness checks, not new confirmatory tests; no multiplicity-adjusted discovery claim is based on them. Code and machine-readable outputs are in \nolinkurl{res_attack/confirmatory_analysis/major_revision/} and \nolinkurl{res_attack/confirmatory_analysis/question_within_deployment/}.

H1 and H2 are the confirmatory tests, and H2 was the single primary contrast named in the preregistration. H3 is reported as a manipulation check. Holm-adjusted budget-specific comparisons are secondary. H4 is reported only as an exploratory directional analysis. The preregistration originally made the confirmatory claim conditional on a directionally consistent critical-budget shift; H4 was reclassified as exploratory after the confirmatory analysis, in response to review, because no observation approached a 50\% crossing and matched coupling estimates are unavailable (Section 3.7). This hierarchy governs the language used below.

\subsection{Deviations from the preregistration}

The study was registered in a repository file rather than with an independent registry, and the file was amended during the study. Table 2 lists every deviation between that original registration and the analysis reported here. The deployment, replacement, and estimator changes were made before the accepted confirmatory data were analyzed, although some followed outcomes from screening or failed runs, as noted. The pooled random-effect change and the H4 reclassification were made after the confirmatory analysis and are therefore post hoc; for transparency the originally specified pooled model is reported alongside the revised one.


\begin{table}[htbp]
\centering
\caption{Deviations from the original preregistration}
\scriptsize
\begin{tabularx}{\textwidth}{>{\raggedright\arraybackslash}X>{\raggedright\arraybackslash}X>{\raggedright\arraybackslash}X>{\raggedright\arraybackslash}X}
\toprule
Item & Original preregistration & Reported study & When and why \\
\midrule
Deployments & Qwen3.5-9B and Llama-3.1-8B via third-party routing; stopping rule: "do not replace it with another model or provider" & Qwen3.5-27B (Alibaba official) and Mistral Small 4 (Mistral official) & 2026-10-28 to 10-29. Llama failed a control-competence gate added on 10-28 after its screening cells were observed; about seven further candidates were screened on six training questions only. A full Qwen3.5-9B confirmatory block on the 18 test questions failed its technical gate (56\% raw-zero spins, 60\% message failures) and was excluded; its per-episode outcomes had been generated. The 27B replacement amendment is dated before any outcome from that deployment. \\
Replacement rule & Pause and diagnose after a gate failure & Recollect the whole block under a dated amendment & Adopted 2026-10-29 and applied to all five failures (Section 3.5). \\
Primary contrast & H2 (downgrade vs none) only & H1 and H2 confirmatory; H3 manipulation check & H1 and H3 were added as co-reported contrasts; Holm correction over H1–H3 leaves all conclusions unchanged (H2 adjusted p unchanged). \\
H2 wording & Downgrade reduces attack success \textit{and} raises the critical budget & Downgrade reduces attack success & The critical-budget half is reported under H4; it is not supported for Mistral (multiplier CI includes 1). \\
H3 wording & Drop is an oracle-like upper bound & Manipulation check & Drop prompts equal the control apart from shuffle order (Section 4.2). \\
H4 & Confirmatory claim conditional on a directionally consistent critical-budget shift in both deployments & Exploratory & Reclassified after the 2026-11-06 analysis, in response to review. Under the original criterion the point estimates would have counted as consistent. \\
Estimator details & Mixed logistic model; terms and random intercepts named & Linear raw budget, control coded \texttt{none}, contrast at 7/3, MAP with stated priors, Wald intervals & Fixed in analysis code before the confirmatory merge, not in the registration text. \\
Pooled random effects & Question random intercept (sharing across deployments not specified) & Question within deployment (primary); shared intercept reported alongside & Changed after review because difficulty profiles are opposite across deployments (Section 3.6). \\
Deployment-specific fits & Not specified & Reported as within-deployment estimates & Added; not preregistered. \\
\bottomrule
\end{tabularx}
\end{table}
\FloatBarrier

\section{Results}

\subsection{Accepted data and technical validity}

All 20 accepted blocks met the prespecified gates, yielding 3,600 episodes. Across 259,200 raw spin calls, 151 returned zero, a rate of 0.058\%. Every raw zero was classified as a request error. There were no exhausted parse failures and no recovered parse retries. Twenty-eight of 57,600 unique-sender message calls failed, a rate of 0.049\%. No valid honest message produced an invalid credential.

\subsection{Descriptive attack-success trajectories}

Under \texttt{none}, observed incorrect-majority rates increased with budget in both deployments (Table 3). Qwen rose from 0.161 in the budget-zero control to 0.194, 0.211, and 0.261 at budgets one, two, and four. Mistral rose from 0.211 to 0.294, 0.311, and 0.361.

Credential-aware trajectories differed by deployment. Qwen downgrade rates were 0.200, 0.194, and 0.189, while its drop rates were 0.178, 0.167, and 0.144. Mistral downgrade rates were 0.228, 0.272, and 0.322, while its drop rates were 0.222, 0.222, and 0.239. These are descriptive proportions; the mixed models provide the prespecified inferential comparisons.


\begin{table}[htbp]
\centering
\caption{Incorrect-terminal-majority rates by deployment, policy, and budget}
\scriptsize
\resizebox{\textwidth}{!}{%
\begin{tabular}{lrrrrrrrrrr}
\toprule
Deployment & Control & None b=1 & None b=2 & None b=4 & Downgrade b=1 & Downgrade b=2 & Downgrade b=4 & Drop b=1 & Drop b=2 & Drop b=4 \\
\midrule
Qwen3.5-27B & 0.161 & 0.194 & 0.211 & 0.261 & 0.200 & 0.194 & 0.189 & 0.178 & 0.167 & 0.144 \\
Mistral Small 4 & 0.211 & 0.294 & 0.311 & 0.361 & 0.228 & 0.272 & 0.322 & 0.222 & 0.222 & 0.239 \\
\bottomrule
\end{tabular}%
}
\end{table}
\FloatBarrier

\texttt{Drop} rates stayed at the control level, as the design implies: drop minus control was +0.017, +0.006, and −0.017 for Qwen and +0.011, +0.011, and +0.028 for Mistral at budgets one, two, and four. This confirms that the attack reached agents only through the routed channel and that removing it restored control-level behavior; it is the manipulation check described in Section 1, not a separate finding about defense strength.

The outcome was highly concentrated by question. For Qwen, 12 of 18 questions never ended in an incorrect majority in any of their 100 episodes, 2 always did, and only 4 varied (\texttt{math\_test\_132}, \texttt{math\_test\_157}, \texttt{math\_test\_21}, \texttt{math\_test\_426}). For Mistral, 8 questions never ended incorrectly, none always did, and 10 varied. The 16–21\% control rates therefore mostly reflect a few questions that were wrong regardless of condition, and the treatment contrasts are carried by the questions whose outcomes varied.

\subsection{Confirmatory tests}

The primary pooled model supports both confirmatory hypotheses (Table 4). Under \texttt{none}, each added external sender multiplied the conditional odds of an incorrect majority by 1.90 (95\% Wald CI 1.59 to 2.27, p = 1.3 x 10\textasciicircum{}-12), supporting H1. At the mean positive budget, \texttt{downgrade} reduced the odds relative to \texttt{none} to 0.48 (95\% CI 0.32 to 0.71, p = 2.3 x 10\textasciicircum{}-4), supporting H2, and the \texttt{drop} manipulation check gave OR 0.20 (0.13 to 0.30). Each deployment shows the same pattern. For Qwen, the budget slope was 3.16 (2.08 to 4.81) and \texttt{downgrade} versus \texttt{none} was 0.27 (0.11 to 0.67, p = .0043); for Mistral, the corresponding estimates were 1.63 (1.34 to 1.99) and 0.56 (0.36 to 0.87, p = .010). Deployment-specific magnitudes depend on the random-effect prior (Section 3.6) and are best read for direction and approximate size.

Within-question randomization tests reach the same conclusion without any model for question effects. Holding the 18 questions, seeds, and budgets fixed, the mean paired risk difference \texttt{none} − \texttt{downgrade} was 0.028 for Qwen (two-sided p = .0017) and 0.048 for Mistral (p = .0043); \texttt{none} − \texttt{drop} was 0.059 and 0.094 (both p < .0001).

The preregistered shared-intercept specification gives the same directions with smaller magnitudes: OR 1.229 (1.119 to 1.349, p = 1.45 x 10\textasciicircum{}-5) for H1, 0.794 (0.639 to 0.986, p = .0373) for H2, and 0.592 (0.474 to 0.740) for the \texttt{drop} check. Because question difficulty runs in opposite directions across the two deployments, a shared intercept pulls these estimates outside the range of both deployment-specific values (Section 3.6), so we treat the question-within-deployment model as primary.


\begin{table}[htbp]
\centering
\caption{Deployment-specific and pooled contrasts}
\scriptsize
\resizebox{\textwidth}{!}{%
\begin{tabular}{lrrrr}
\toprule
Scope & Contrast & OR & 95\% CI & p \\
\midrule
Qwen3.5-27B & H1: budget slope under none & 3.16 & 2.08 to 4.81 & < .0001 \\
Qwen3.5-27B & H2: downgrade vs none & 0.27 & 0.11 to 0.67 & .0043 \\
Qwen3.5-27B & H3 (check): drop vs none & 0.06 & 0.02 to 0.15 & < .0001 \\
Mistral Small 4 & H1: budget slope under none & 1.63 & 1.34 to 1.99 & < .0001 \\
Mistral Small 4 & H2: downgrade vs none & 0.56 & 0.36 to 0.87 & .010 \\
Mistral Small 4 & H3 (check): drop vs none & 0.28 & 0.18 to 0.45 & < .0001 \\
Pooled, question within deployment & H1: budget slope under none & 1.90 & 1.59 to 2.27 & 1.3 x 10\textasciicircum{}-12 \\
Pooled, question within deployment & H2: downgrade vs none & 0.48 & 0.32 to 0.71 & 2.3 x 10\textasciicircum{}-4 \\
Pooled, question within deployment & H3 (check): drop vs none & 0.20 & 0.13 to 0.30 & 2.8 x 10\textasciicircum{}-14 \\
Pooled, preregistered shared question & H1: budget slope under none & 1.229 & 1.119 to 1.349 & 1.45 x 10\textasciicircum{}-5 \\
Pooled, preregistered shared question & H2: downgrade vs none & 0.794 & 0.639 to 0.986 & .0373 \\
Pooled, preregistered shared question & H3 (check): drop vs none & 0.592 & 0.474 to 0.740 & 4.10 x 10\textasciicircum{}-6 \\
\bottomrule
\end{tabular}%
}
\end{table}
\FloatBarrier

\subsection{Secondary policy contrasts}

In the question-within-deployment model, after Holm correction across the nine secondary comparisons, \texttt{drop} differed from \texttt{none} at every budget (Holm-adjusted p = .028, 2.4 x 10\textasciicircum{}-10, and 8.7 x 10\textasciicircum{}-14 at budgets one, two, and four), and \texttt{downgrade} differed from \texttt{none} at budgets two (OR 0.54, Holm-adjusted p = .0070) and four (OR 0.27, p = .00054) but not one (OR 0.76, p = .36). \texttt{Drop} also differed from \texttt{downgrade} at budgets two and four. In the preregistered shared-intercept model, only \texttt{drop} versus \texttt{none} at budgets two and four passed the corrected threshold (Holm-adjusted p = .000635 and 2.05 x 10\textasciicircum{}-5); \texttt{downgrade} versus \texttt{none} at budget four did not (Holm-adjusted p = .150).

These comparisons are points on the fitted linear budget-by-defense line, not tests of individual cells, so the value at budget four is constrained by the linear trend. The categorical GEE in Section 4.5 provides a cell-based check.

\subsection{Robustness and influence diagnostics}

\textbf{Question-level influence.} The Wald p values in Table 4 assume a treatment effect common across questions, and Qwen's effects are carried by its four informative questions, whose budget slopes differ. With question-clustered GEE inference the Qwen contrasts are not individually significant (p ≈ .09 to .13), and Mistral's H2 is borderline (p ≈ .07). Omitting each pair of Qwen's informative questions in turn, H1 and the \texttt{drop} check remained significant in five of six refits and H2 in three of six; omitting \texttt{math\_test\_157} and \texttt{math\_test\_21} together removed all three effects. The Qwen results therefore establish the effects for this question set, and a larger question set is needed for population-level precision.

Prior choice did not change any direction or significance within the tested grid (Table 5). In the question-within-deployment model, H1 ORs ranged from 1.76 to 1.91, H2 from 0.48 to 0.52, and H3 from 0.20 to 0.23 across the nine-point grid, with no optimizer warnings; the shared-intercept model ranged even less. As Section 3.6 notes, this stability holds only within prior scales that keep the joint MAP well defined, and the random-effect SDs themselves move with the prior.

The categorical-cell GEE removed the linear-budget assumption and used sandwich covariance over 18 question clusters. Under \texttt{none}, the ORs relative to control were 1.423 at budget one (95\% CI 1.077 to 1.882), 1.557 at budget two (1.070 to 2.267), and 1.998 at budget four (1.244 to 3.209). Averaging the three positive-budget log-odds differences gave OR 0.817 for downgrade versus none (0.701 to 0.952, p = .00966) and OR 0.647 for drop versus none (0.484 to 0.866, p = .00336). With a Mancl–DeRouen small-sample correction and a t(17) reference, these remain significant (H2 p = .026; drop p = .013). The GEE is a marginal estimand and shares one working correlation across deployments, so its magnitudes are not directly comparable with the conditional ORs.

Leave-one-question-out fits of the question-within-deployment model preserved all directions and significance: H1 ORs ranged from 1.74 to 2.18, H2 from 0.39 to 0.56 (largest p = .005), and H3 from 0.15 to 0.24. Because four questions never produced an incorrect majority in either deployment, omitting any of them reproduces the full fit, so the 18 refits contain about 14 distinct influence checks. In the preregistered shared-intercept model, H2's p value exceeded .05 after four of the 18 omissions (range .025 to .126). That fragility is largely a consequence of the shared intercept rather than of H2 itself.

Substituting the five failed first attempts for their accepted replacements weakened but did not reverse the results: in the question-within-deployment model, H1 OR 1.59, H2 OR 0.62 (p = .011), and H3 OR 0.28; in the shared-intercept model, H2 OR 0.84 (p = .124). Removing only the most contaminated cells of the failed attempts did not restore the accepted values, so the attenuation looks like sampling variation between equally valid draws rather than a contamination artifact. H2's precision therefore depends on which draw of those five seeds is analyzed.


\begin{table}[htbp]
\centering
\caption{Post-review robustness analyses}
\scriptsize
\resizebox{\textwidth}{!}{%
\begin{tabular}{lrrr}
\toprule
Analysis & H1 & H2 & H3 (check) \\
\midrule
Prior grid, question within deployment, OR range & 1.76–1.91 & 0.48–0.52 & 0.20–0.23 \\
Prior grid, shared question, OR range & 1.224–1.229 & 0.793–0.796 & 0.592–0.595 \\
Question-cluster GEE & 1.423, 1.557, 1.998 at b=1,2,4 vs control & 0.817 [0.701, 0.952] & 0.647 [0.484, 0.866] \\
Leave-one-question-out, question within deployment, OR range & 1.74–2.18 & 0.39–0.56 & 0.15–0.24 \\
Leave-one-question-out, shared question, OR range & 1.186–1.268 & 0.767–0.840 (p .025–.126) & 0.544–0.640 \\
Failed attempts substituted, question within deployment & 1.59 & 0.62 (p = .011) & 0.28 \\
Failed attempts substituted, shared question & 1.18 & 0.84 (p = .124) & 0.63 \\
\bottomrule
\end{tabular}%
}
\end{table}
\FloatBarrier

\subsection{Exploratory critical budgets and H4}

Figure 2 shows the observed proportions and fitted nondecreasing curves. Every estimated 50\% crossing lies beyond the maximum tested budget of four, and the highest fitted rate at budget four was 0.37. Only three curves had slopes distinguishable from zero: Qwen \texttt{none} (p = .018), Mistral \texttt{none} (p = .003), and Mistral \texttt{downgrade} (p = .009). The unconstrained slopes for Qwen \texttt{downgrade} (p = .63), Qwen \texttt{drop} (negative, p = .55), and Mistral \texttt{drop} (p = .54) were not distinguishable from zero. For those curves the critical budget is either infinite by construction or rests on a near-zero slope, so we report them as "no detectable positive slope" rather than as numbers. For Qwen, the \texttt{none} critical budget was 11.10 (bootstrap 95\% CI 4.72 to infinity). For Mistral, the estimates were 7.18 under \texttt{none} (2.57 to 24.98) and 8.93 under \texttt{downgrade} (4.07 to 23.73), a downgrade-to-none multiplier of 1.24 (0.60 to 2.18).

\begin{figure}[htbp]
\centering
\includegraphics[width=\textwidth]{../res_attack/confirmatory_analysis/confirmatory_curves.png}
\caption{Observed attack success and fitted nondecreasing curves. Points show incorrect-terminal-majority rates at budgets 0, 1, 2, and 4. Lines are fitted curves. The horizontal reference is 50\% attack success; the vertical reference marks the largest tested budget. All crossings to the right of budget four are extrapolations.}
\end{figure}
\FloatBarrier


\begin{table}[htbp]
\centering
\caption{Exploratory extrapolative 50\% critical budgets}
\scriptsize
\resizebox{\textwidth}{!}{%
\begin{tabular}{lrrrr}
\toprule
Deployment & None & Downgrade & Drop & Downgrade / none \\
\midrule
Qwen3.5-27B & 11.10 [4.72, infinity] & no detectable positive slope & no detectable positive slope & not estimable \\
Mistral Small 4 & 7.18 [2.57, 24.98] & 8.93 [4.07, 23.73] & no detectable positive slope & 1.24 [0.60, 2.18] \\
\bottomrule
\end{tabular}%
}
\end{table}
\FloatBarrier

For Mistral, the downgrade-to-none multiplier exceeds one, but its interval includes one. For Qwen, the defenses flatten the curve so that no positive slope is detectable, which implies an upward crossing but adds nothing beyond H2. The evidence does not establish quantitative calibration: all crossings are outside the observed range, the bootstrap percentile intervals are unreliable when the slope is near zero, and the available coupling file pertains to Qwen3.5-9B rather than Qwen3.5-27B and contains no Mistral estimate.

\section{Discussion}

\subsection{Principal findings}

Two findings answer the research question, and a manipulation check supports their interpretation. First, the repetition-and-reach budget matters when a router trusts self-declared relationships: each one-unit increase raised incorrect-majority odds under \texttt{none} in both deployments. Second, \texttt{downgrade}, which keeps the content but removes the claimed relationship and labels it as unverified, reduced that effect in both deployments. Third, \texttt{drop} returned attack success to control levels, confirming that the attack operated only through the routed channel. Mistral's critical-budget point estimate moved upward under \texttt{downgrade}, but its interval includes one and all crossings are extrapolations.

The result is best understood as a routing effect. The policies do not improve the base model, identify false mathematics, or train agents to resist persuasion. They change the context in which a message reaches the model. Under \texttt{none}, an external sender selects the positive-relation channel by assertion. Under \texttt{downgrade}, the content remains visible but loses the asserted relationship and gains a provenance notice. Under \texttt{drop}, the content never enters deliberation. The routing decision itself is deterministic and needs no semantic detector. The \texttt{downgrade} effect, however, is realized through how the model reads the relabeled prompt, so it is partly a prompt-level effect (Section 5.3).

\subsection{What the budget manipulates: repetition and reach, not identities}

Multi-agent reasoning often benefits from multiple traces because agents may make different errors \cite{du2024debate}, \cite{liang2024divergent}. That benefit depends on diversity and partial independence. Strong prompting can already narrow the gap between one agent and a discussion \cite{wang2024bounds}, and homogeneous interaction can reinforce a shared mistake \cite{oh2026belief}. A system should therefore avoid equating message count with evidentiary diversity.

The treatment targets that assumption by adding identical copies of the same false option. When the router accepted each relationship claim, those copies entered the positive-relation prompt section, and more senders also reached more honest agents. Because the prompt shows no sender identifiers, the H1 effect cannot reflect a perceived number of distinct speakers. It is consistent with evidence that repeating a wrong answer, even without naming a speaker, produces most of the conformity seen in peer-pressure benchmarks \cite{hu2026speaker}, and with group-conformity \cite{choi2025conformity} and committed-minority \cite{ashery2025conventions} results, but consistency is not mechanism identification. In particular, the design does not separate repetition within a receiver's prompt from reach across receivers.

In an operational system, this bundle is a plausible abuse path: one principal multiplies identifiers to multiply copies of its message. For science, the distinction matters. An identity-specific claim would require prompts that render sender identifiers, a fixed-total-volume design that varies the number of distinct identifiers, and a fixed-identifier design that varies repetition.

\subsection{Downgrade and drop encode different policies}

Downgrade retains content from an unknown relationship. This can be appropriate in open systems where a new sender may contribute useful information but should not inherit the influence associated with an enrolled edge. It separates content availability from relationship privilege. The protective H2 estimate shows that the full downgrade presentation reduces the attack even when the message remains accessible. It does not show which part of that presentation matters. Downgrade simultaneously removes the positive-relationship framing, adds the notice that the messages "failed relationship verification," segregates them into their own section, and moves them next to the problem statement. Each component could plausibly contribute, and the position change could plausibly work in either direction. The policy is therefore closer to a provenance-labeling defense such as spotlighting \cite{hines2024spotlighting} than to a pure routing change. A neutral-header ablation in the same position, and a variant that keeps the warning but places the messages in the original section, would separate these components. Prompt Infection's finding that source tagging alone gave little protection \cite{lee2024promptinfection} suggests that the explicit verification notice may matter more than the label change.

Drop enforces a closed-membership rule. In this design, \texttt{drop} reproduces the control condition, so its effect equals the full attack effect and serves as an upper bound on what any admission policy could achieve against this traffic. A deployment with strict authorization requirements may prefer this rule. An open collaboration platform may not, because valid but previously unseen contributors would also be excluded unless an onboarding mechanism issued credentials. In open ecosystems, such as open-source agent swarms, \texttt{drop} is therefore likely to be too costly, because it shuts out every new benign participant; \texttt{downgrade} offers a form of graceful degradation, keeping new information available while withholding relationship privilege, and H2 shows that this still removes much of the attack's effect.

No universal ranking follows from the current data. The mathematics corpus contains no benign unknown senders, and the outcome does not measure helpful information lost under drop. A complete policy evaluation should cross malicious and benign unknown-origin traffic, measure task utility and security jointly, and vary the cost of credential acquisition. Until then, drop can be described as more protective against this attack, not as categorically superior.

\subsection{Cross-deployment replication and heterogeneity}

H1, H2, and the \texttt{drop} check have the same direction in Qwen3.5-27B and Mistral Small 4. This replication matters because API-based agent experiments can depend on endpoint behavior, generation settings, and provider infrastructure. It is, however, a replication on the same 18 questions, and for Qwen the information comes from four of them. The accepted-block gates also show that technical reliability was high in both collections after complete replacement of failed blocks.

The magnitudes differ substantially. Qwen's deployment-specific budget slope and policy contrasts are larger, and its observed downgrade trajectory is nearly flat. Mistral's downgrade trajectory rises with budget while remaining below \texttt{none}. Several explanations are possible: differences in model behavior, system prompts as processed by the provider, sampling implementation, or serving-stack details. The design cannot separate them. Calling one model family more vulnerable or one defense more effective for that family would therefore be a causal overreach.

\subsection{Architectural implications}

The main architectural implication is to treat identity and relationship as authorization data. A natural-language statement such as “I am your trusted collaborator” should not create a privileged edge. The sender identity, receiver identity, authorized relationship, and freshness value should be covered by a verifiable object whose validation occurs outside the LLM. This follows the logic of authenticated delegation \cite{south2025delegation}, verifiable agent credentials \cite{garzon2026credentials}, standardized message authentication \cite{nist2008hmac}, and classical membership assumptions \cite{douceur2002sybil}, \cite{castro1999pbft}.

The following design arguments motivate the token but were not tested; in this experiment a membership check would have routed traffic identically. Relationship claims deserve separate protection from sender identity. Authenticating only a sender can show that a known principal transmitted a message, but not that the principal holds the asserted role for a particular receiver. A token that binds the edge, rather than only the node, addresses this. The distinction becomes important in systems where one agent may supervise, oppose, audit, or merely observe different peers, and where an enrolled agent might try to upgrade its relationship to a particular receiver. Testing it requires an attacker that presents a valid sender identity with a missing, wrong-edge, or substituted-relation token.

The control is compatible with defense in depth. Authenticated routing can operate alongside prompt-injection isolation, content scanning, rate limits, anomaly detection, and Byzantine-aware or credibility-weighted aggregation \cite{zheng2025byzantine}, \cite{ebrahimi2025credibility}. Admission control does not stop influence that has already been admitted: under \texttt{none} and \texttt{downgrade}, step-two honest messages were generated from inboxes that contained external content, so some attacker influence reached receivers through validly credentialed honest relays. The measured effects include this propagation. Its contribution is to preserve provenance before those later mechanisms run. If one external principal creates many labels, identity enrollment and rate policy can also aggregate them at a principal level. If an enrolled principal becomes malicious, content and consensus defenses remain necessary because authentication will correctly report who sent the harmful message rather than reject it.

\section{Limitations and Future Work}

The strongest construct-validity limitation is that the credential mechanism was never challenged. Every external envelope was untokened and came from a non-member identifier, so a membership check would have produced identical routing; receiver binding, relation binding, and freshness were not tested, and the implemented verifier did not enforce freshness. The budget manipulation is also bundled: the prompt shows no sender identifiers, and each added sender added identical copies and reached more receivers, so the effect reflects repetition and reach and does not separate the two. The \texttt{downgrade} presentation bundles label removal, a verification notice, segregation, and position. The corpus contains only 18 binary-answer, text-only mathematics questions, and agents were instructed to give quick judgments rather than reason, which likely increases sensitivity to social cues. Open-ended tasks, subjective judgments, tool use, deliberate reasoning, and long-horizon planning may respond differently. Each episode used eight honest agents, a fixed graph density, two synchronous updates, and one standardized false-option message. Larger or dynamic networks could change both attack reach and defense behavior.

All tested external-identity budgets were at most four. Every 50\% critical budget is therefore model-based extrapolation, and several curves have no detectable positive slope. H4 cannot support numerical calibration. The HMAC authority is an experimental implementation with a reproducibly derived key, not a production identity system. The study does not cover stolen keys, compromised enrolled agents, credential sharing, replay, revocation, colluding authenticated nodes, impersonation embedded inside credentialed messages \cite{cui2025madspear}, or malicious content sent with valid authorization.

The MAP/Laplace intervals depend on the specified random-effects model and use approximate Wald inference. Most questions produced the same outcome in every episode, so random-effect scales and conditional OR magnitudes depend on the prior. Only 4 Qwen and 10 Mistral questions were informative, and with question-clustered inference the Qwen contrasts are not individually significant. Within-question randomization tests support the effects for this question set, but generalization to new questions rests on few informative items, especially for Qwen. The pooled random-effect structure was revised after review, and H2's precision depends on which attempt of the five replaced blocks is analyzed. The original sensitivity calculation used a two-independent-proportion approximation rather than a simulation of the full crossed mixed-effects design.

The two deployments differ in model family, provider, and serving stack. Their estimates are independent replications on the same questions, not a controlled model-family comparison. The deployments themselves differ from those in the original preregistration (Section 3.7). Five model-seed blocks required complete recollection after technical failure. Whole-block replacement prevented selective cell reuse, but the amendments were repository records rather than an independently timestamped preregistration. Cells within each block were run in a fixed order (control, then budgets one, two, and four), so budget is correlated with collection time; request errors clustered in late cells, although they stayed far below the gates.

Finally, the outcome suite emphasizes attack resistance. It does not quantify the benign utility of accepting an unknown sender, so it cannot establish a general preference between downgrade and drop. Future work should test an attacker that spoofs enrolled identities or presents wrong-edge or replayed tokens; add a neutral-header downgrade arm; render sender identifiers and vary them at fixed message volume; test benign unknown senders alongside malicious ones; add larger budgets, adaptive targeting, varied rhetoric, credential compromise, dynamic graphs, longer horizons, and a larger or harder question set; randomize cell order; and compare models behind a common serving stack. A matched physical-model study should estimate coupling for each live deployment and sample budgets around the predicted transition.

\section{Conclusion}

This paper examined whether a router should accept an external message's self-declared relationship. Across 3,600 accepted episodes in two official LLM deployments, adding untokened external senders—which, because sender identifiers are not shown to the model, meant adding identical false-message copies and reaching more agents—raised incorrect-majority odds when the claimed relationship was trusted. Downgrading that traffic to a labeled unverified-source channel reduced the effect in both deployments, and dropping it restored control-level behavior. These conclusions hold under within-question randomization tests and across the robustness analyses, but for Qwen they rest on four informative questions.

The practical insight is architectural. Enrollment, authentication, relationship authorization, routing, and semantic judgment should remain separate, and a natural-language or metadata claim should not by itself place a message in a trusted relationship channel. In this experiment that rule was exercised only against non-members, where membership gating and credential checking coincide. Whether binding credentials to receivers, relations, and freshness adds protection against enrolled or impersonating senders remains to be tested.

The scope is deliberately limited. The budget bundled repetition and reach, the downgrade arm bundled label removal with a verification notice, benign unverified traffic was not tested, the question set is small, and every fitted 50\% critical budget lies outside the observed range. The evidence therefore does not identify an identity mechanism, establish a security-utility optimum, demonstrate universal immunity, validate physical calibration, or rank model families. Within the tested attack bundle, refusing to honor self-declared relationships from non-members before prompt construction provides a reproducible way to reduce unauthenticated external influence. The next decisive experiments should challenge the credential with spoofed and replayed tokens, isolate the components of \texttt{downgrade}, and cross malicious with benign unknown-origin traffic, thereby closing the gap between a metadata-trusting baseline and a fully authenticated, open agent ecosystem.

\section{Questions for Feedback}

This draft is circulated for comments. The points on which advice would help most are:

1. \textbf{Contribution and framing.} Is "refusing self-declared relationships from non-members before prompting," measured on collective outcomes, a sufficient contribution for a workshop paper, or should the credential-challenge experiment (spoofed members, wrong-edge and replayed tokens) come first?
2. \textbf{Next experiment.} Which should take priority: the neutral-header \texttt{downgrade} ablation, the credential-challenge attacker, benign unknown senders to measure utility, or a larger and harder question set to strengthen inference beyond Qwen's four informative questions?
3. \textbf{Statistical presentation.} Should the paper lead with the revised pooled model (question intercepts within deployment), or with the deployment-specific estimates and within-question randomization tests, keeping the pooled models secondary?
4. \textbf{Venue.} Does the work fit better at a security workshop or an agents/ML workshop, and is the link to the \textit{Physics of Agents} coupling framework worth developing into a quantitative test?
5. \textbf{Transparency.} Is the deviations table (Section 3.7) sufficient for a study that changed deployments after registration?

\section{Data and Code Availability}

The accepted-block manifest, SHA-256 source hashes, merged episode table, analysis scripts, robustness outputs, bootstrap draws, tables, prompt appendix, and figure sources are provided at \nolinkurl{https://github.com/Ben-Tseng/physics-of-agents}, a fork of the \textit{Physics of Agents} repository \cite{el2026physics}. The archival revision will be tagged before submission. Post-review analyses are produced by \nolinkurl{res_attack/question_within_deployment_refit.py} and \nolinkurl{res_attack/replacement_sensitivity.py}. The analysis enumerates the 20 accepted blocks explicitly and excludes all predecessor directories replaced by full-block recollection.

\section{Acknowledgments}

OpenAI Codex and Anthropic Claude assisted with code inspection, numerical cross-checking, literature organization, drafting, and pre-submission review. The author designed the study, controlled data collection, verified all analyses, numbers, and references, and takes full responsibility for the content.

\begingroup
\small
\bibliographystyle{IEEEtran}
\bibliography{references}
\endgroup
\end{document}
